\section{Complete finite-graph specification}
\label{app:certificate}

This appendix gives the full finite arithmetic specification for
Propositions~\ref{prop:large-computation} and \ref{prop:finite}.
The pseudocode uses exact integers, finite sets, dictionaries and
lexicographic order. Its half-open \texttt{range(a,b)} contains
$a,a+1,\ldots,b-1$; \texttt{product} is Cartesian product.
There is no imposed orbit-length bound. A word is canonicalized by
rotation only.

\subsection{Cycle extraction and its correctness}

For an ordered finite symbol set $S$ and values $V[y]$, the graph
has the edge $(x,y)\mapsto(y,V[y]-x)$ when the last coordinate
belongs to $S$. The symbols and values are either integers or pairs
of integer coefficients. Subtraction is ordinary subtraction in the
first case and componentwise subtraction in the second.

\begin{Verbatim}[fontsize=\footnotesize]
from itertools import product

def minus(u, v):
    if isinstance(u, tuple):
        return (u[0] - v[0], u[1] - v[1])
    return u - v

def cycles(S, V):
    S = tuple(sorted(S))
    allowed = set(S)
    seen = set()
    answer = []
    for start in product(S, repeat=2):
        if start in seen:
            continue
        point = start
        path = []
        position = {}
        while point not in seen:
            if point in position:
                orbit = path[position[point]:]
                word = tuple(x for (x, y) in orbit)
                word = min(word[i:] + word[:i]
                           for i in range(len(word)))
                answer.append(word)
                break
            position[point] = len(path)
            path.append(point)
            x, y = point
            z = minus(V[y], x)
            if z not in allowed:
                break
            point = (y, z)
        seen.update(path)
    return sorted(answer, key=lambda w: (len(w), w))
\end{Verbatim}

\begin{lemma}[Exact extraction]
\label{lem:extraction}
The procedure \texttt{cycles} terminates and returns every directed
cycle of the partial graph once, and no other words. Each returned
word has length equal to its least ordinary period.
\end{lemma}
\begin{proof}
The graph has finitely many vertices. Within one traversal, either
an image leaves the graph, the path reaches an already completed
traversal, or a vertex repeats on the current path; one of these
events occurs after finitely many steps. A repeat on the current
path closes exactly its segment from the first occurrence to the
repeat. Its vertices are distinct before closure, so it is one
whole cycle with its least period.

Consider any cycle and the first traversal that reaches it.
No vertex of that cycle is yet in a completed traversal.
Following the deterministic edges therefore closes and records it.
A subsequent path reaching that cycle stops at a completed vertex
and cannot record it again. Distinct cycles cannot intersect in a
functional graph. Marking every traversed vertex completed thus
loses no unrecorded cycle and creates no duplicate. The outer loop
visits every vertex, so every cycle is eventually reached.
Canonical rotation changes only its chosen starting coordinate,
not its orientation or its length.
\end{proof}

\subsection{Small-diameter enumeration}

The following code computes the input counts, the endpoint-filtered
point maxima, period sets, per-cycle symbol maxima, and every global
maximizer. The value eleven is not supplied as an input or assertion.

\begin{Verbatim}[fontsize=\footnotesize]
def small_certificate():
    rows = []
    global_maximum = -1
    equality = []
    all_periods = set()
    for D in range(2, 16):
        cases = kept = 0
        maximum = 0
        max_symbols = 0
        periods = set()
        for r, A, q in product(range(-3, D + 4),
                               range(2 * D + 1),
                               range(-2, 3)):
            if not (0 <= A + q * D <= 2 * D):
                continue
            cases += 1
            S = [t for t in range(D + 1)
                 if abs(t * (t-D) * (t-r)) <= 2 * D]
            V = {t: A + q*t + t*(t-D)*(t-r) for t in S}
            C = cycles(S, V)
            used = {t for w in C for t in w}
            if not (0 in used and D in used):
                continue
            kept += 1
            count = sum(len(w) for w in C)
            maximum = max(maximum, count)
            for w in C:
                periods.add(len(w))
                all_periods.add(len(w))
                max_symbols = max(max_symbols, len(set(w)))
            if count > global_maximum:
                global_maximum = count
                equality = []
            if count == global_maximum:
                equality.append((D, r, A, q, C))
        rows.append((D, cases, kept, maximum,
                     sorted(periods), max_symbols))
    return rows, global_maximum, equality, sorted(all_periods)
\end{Verbatim}

The ranges are exactly \eqref{eq:finite-ranges}; the symbol filter
is \eqref{eq:small-symbols}. A retained case has at least one cycle,
so every per-cycle maximum is defined. The output rows are
Table~\ref{tab:small}; the global maximizer and its words are
\eqref{eq:max-tuple}--\eqref{eq:max-words}. The period union is
$\{1,2,3,4,6\}$ and every row has symbol maximum three.

\subsection{Symbolic large-diameter enumeration}

Represent $kD+s$ by $(k,s)$. The symbol $D$ is $(1,0)$ and an
integer symbol $t\in T_r$ is $(0,t)$. Formula
\eqref{eq:symbolic-value} defines its image coefficients.
The extra statistic \texttt{error} measures the maximum absolute
constant discrepancy among all edge tests. This statistic checks
the arithmetic bound, but the independent proof of uniformity
remains Lemma~\ref{lem:no-carry}.

\begin{Verbatim}[fontsize=\footnotesize]
def symbolic_certificate():
    rows = []
    endpoint_patterns = []
    for r in (-1, 0, 1, 2, 3):
        T = [t for t in range(4)
             if t == 0 or abs(t * (t-r)) <= 2]
        S = [(0, t) for t in T] + [(1, 0)]
        Sigma = {(u[0] + v[0], u[1] + v[1])
                 for u, v in product(S, repeat=2)}
        cases = kept = maximum = max_symbols = error = 0
        periods = set()
        for (k, s), q in product(sorted(Sigma), range(-2, 3)):
            if (k + q, s) not in Sigma:
                continue
            cases += 1
            V = {(0, t): (k - t*(t-r),
                          s + q*t + t*t*(t-r)) for t in T}
            V[(1, 0)] = (k + q, s)
            for y, u, v in product(S, repeat=3):
                error = max(error, abs(V[y][1]-u[1]-v[1]))
            C = cycles(S, V)
            count = sum(len(w) for w in C)
            maximum = max(maximum, count)
            for w in C:
                periods.add(len(w))
                max_symbols = max(max_symbols, len(set(w)))
            used = {t for w in C for t in w}
            if (0, 0) in used and (1, 0) in used:
                kept += 1
                endpoint_patterns.append((r, (k, s), q, C))
        rows.append((r, T, cases, kept, maximum,
                     sorted(periods), max_symbols, error))
    return rows, endpoint_patterns
\end{Verbatim}

The five values of \texttt{error} are $4,3,8,8,14$, all less than
sixteen. They are no larger than the deliberately looser analytic
bounds $6,5,12,8,14$. The other row entries give
Table~\ref{tab:large}. Decoding $(0,t)$ as $t$ and $(1,0)$ as
$D$, and changing only cyclic starting positions, gives exactly
the 29 cases summarized in Table~\ref{tab:large-patterns}.

Both enumeration procedures call the same completely specified
cycle routine. Every vertex, parameter tuple and comparison is
finite and explicit. Sections~\ref{sec:secant}--\ref{sec:finite}
prove why no original polynomial escapes their ranges or the
separate affine cases. This is the required distinction between
an exhaustive finite certificate and a coefficient experiment.
